\appendix

\section{Useful tables}\label{app:useful-tables}

This appendix collects the quantities and formulae that recur across the exercise
sheets. Particle masses, widths and branching ratios are \emph{not} reproduced here:
they are maintained by the Particle Data Group and should be looked up at
\url{http://pdglive.lbl.gov}. What follows is the derived material---constituent masses,
naming conventions, kinematic formulae---that is scattered through the chapters and is
tedious to hunt down mid-calculation.

% ------------------------------------------------------------
\subsection{Constituent quark masses}
\label{app:quark-masses}

The constituent masses are effective masses acquired by quarks inside hadrons, and are
much larger than the bare (current) masses generated by the Higgs mechanism. They are the
masses to use in quark-model estimates of magnetic moments and hyperfine splittings.

The current and constituent masses are tabulated in Table~\ref{tab:quark-masses}.
For reference, the values extracted from the baryon magnetic moments in
Sect.~\ref{subsec:magnetic-moments} are $m_u = m_d \approx 327$~MeV and
$m_s \approx 510$~MeV, consistent with the $\sim 350$~MeV and $\sim 500$~MeV
quoted there.

% ------------------------------------------------------------
\subsection{Relativistic kinematics}
\label{app:kinematics}

Collected from Sect.~\ref{subsec:rel-kinematics}. Natural units throughout,
$\hbar = c = 1$. The invariant variables are treated by Halzen and Martin
\cite[Ch.~4]{HalzenMartin1984}, the experimental relations by Perkins
\cite[Ch.~2]{Perkins2000}, and the invariant-mass relations behind the Dalitz entries by
Amsler \cite[App.~B]{Amsler2018}.

\begin{table}[htbp]
    \centering
    \caption{Kinematic relations used throughout the sheets.}
    \label{tab:kinematics}
    \renewcommand{\arraystretch}{1.5}
    \begin{tabular}{ll}
        \hline
        Quantity & Expression \\
        \hline
        Mandelstam invariants
            & $s=(P_1+P_2)^2$, $t=(P_1-P_3)^2$, $u=(P_1-P_4)^2$ \\
        Sum rule
            & $s+t+u=\sum_i m_i^2$ \\
        Virtuality
            & $s$-channel $q^2=s>0$ (timelike); $t$-channel $q^2=t<0$ (spacelike) \\
        Fixed target
            & $s = m_1^2 + m_2^2 + 2 m_2 E_1$ \\
        Real photon on proton
            & $s = M_p^2 + 2 M_p E_\gamma$ \\
        Threshold beam energy
            & $E_1^{\mathrm{thr}} = \dfrac{\left(\sum_f m_f\right)^2 - m_1^2 - m_2^2}{2m_2}$ \\
        Two-body decay energies
            & $E_1 = \dfrac{M^2+m_1^2-m_2^2}{2M}$ \\
        Breakup momentum
            & $|\vec q\,| = \dfrac{\sqrt{[M^2-(m_1+m_2)^2][M^2-(m_1-m_2)^2]}}{2M}$ \\
        Radiative photon energy
            & $E_\gamma = \dfrac{m_i^2-m_f^2}{2m_i}$ \\
        Decay length
            & $d = \beta\gamma\, c\tau = \dfrac{p}{m}c\tau$;\quad
              $N(\ell)=N_0\,e^{-\ell/\beta\gamma c\tau}$ \\
        Track curvature
            & $p\,[\mathrm{GeV}] = 0.3\, B\,[\mathrm{T}]\, R\,[\mathrm{m}]$ \\
        Drell--Yan pair mass
            & $M^2 = x_1 x_2 s$ \\
        \hline
    \end{tabular}
\end{table}

% ------------------------------------------------------------
\subsection{Dalitz plot formulae}
\label{app:dalitz-formulae}

For a parent of mass $M$ decaying to three bodies of masses $m_1$, $m_2$, $m_3$
(Sect.~\ref{subsubsec:dalitz-kinematics}):
\begin{align}
    m_{12}^2 + m_{13}^2 + m_{23}^2 &= M^2 + m_1^2 + m_2^2 + m_3^2, \notag\\
    m_{ij}^2\big|_{\max} &= (M-m_k)^2, \qquad
    m_{ij}^2\big|_{\min} = (m_i+m_j)^2, \notag\\
    m_{23}^2\big|_{\substack{\max\\\min}}
    &= \left(E_2^*+E_3^*\right)^2
       - \left(|\vec p_2^{\,*}| \mp |\vec p_3^{\,*}|\right)^2 .
\end{align}
The sum rule makes the third invariant mass a diagonal axis on a plot of $m_{12}^2$
against $m_{23}^2$. The boundary corresponds to collinear momenta; the interior to
configurations in which the three momenta close into a triangle.

% ------------------------------------------------------------
\subsection{Baryon nomenclature}
\label{app:baryon-naming}

Baryons are named by isospin and strangeness (Sect.~\ref{subsec:baryon-naming}):
$N$ ($I=\tfrac12$, $S=0$), $\Delta$ ($I=\tfrac32$, $S=0$), $\Lambda$ ($I=0$, $S=-1$),
$\Sigma$ ($I=1$, $S=-1$), $\Xi$ ($I=\tfrac12$, $S=-2$), $\Omega$ ($I=0$, $S=-3$).

The older $\pi N$ scattering notation $L_{2I\,2J}(\text{mass})$ uses spectroscopic
letters for $L$ and \emph{twice} the isospin and total angular momentum as subscripts;
the conversion to modern names is tabulated in Table~\ref{tab:l2i2j}. Parity follows from
$P=(-1)^{L+1}$, so $S$ and $D$ waves are negative-parity and $P$ and $F$ waves positive.

Clebsch--Gordan coefficients for the couplings needed on the sheets
($\tfrac12\otimes\tfrac12$, $1\otimes\tfrac12$, $\tfrac32\otimes\tfrac12$,
$1\otimes 1$, $\tfrac32\otimes 1$, $2\otimes\tfrac12$) are reproduced in
Fig.~\ref{fig:cg_table}.

% ------------------------------------------------------------
\subsection{Exercise sheet cross-reference}
\label{app:sheet-map}

Where to look for the material behind each exercise. Sheet numbers refer to the
SS~2026 problem sets.

\begin{longtable}{llL{7.2cm}}
    \hline
    Sheet & Topic & Sections \\
    \hline
    \endhead
    1.1 & Standard Model particles      & Sect.~\ref{subsec:history-quark-model} \\
    1.2 & Spin addition                 & Sect.~\ref{subsec:addition} \\
    1.3 & Yukawa range                  & Sect.~\ref{subsec:yukawa} \\
    1.4 & Feynman diagrams, $s$/$t$ channel & Sect.~\ref{subsubsec:mandelstam} \\
    \hline
    2.1 & Cross section, luminosity     & Sect.~\ref{subsec:scattering-reminder} \\
    2.2 & Invariant mass from $2\gamma$ & Sect.~\ref{subsec:invariant-mass} \\
    2.3 & Bubble chamber, curvature     & Sect.~\ref{subsec:bubble-chamber} \\
    2.4 & Parity of the charged pion    & Sect.~\ref{subsec:pion-parity} \\
    2.5 & Missing mass                  & Sect.~\ref{subsec:invariant-mass} \\
    \hline
    3.1 & Isospin, $\pi N$ system       & Sect.~\ref{subsec:su2},
                                          Eq.~\eqref{eq:piN-decomposition} \\
    3.2 & Strangeness production, thresholds & Sect.~\ref{subsubsec:thresholds},
                                          Sect.~\ref{subsec:accelerator-discovery} \\
    3.3 & Young tableaux                & Sect.~\ref{subsec:youngtableaux} \\
    \hline
    4.1 & Gell-Mann--Okubo              & Sect.~\ref{subsec:gmo},
                                          Sect.~\ref{subsec:gmo-mesons} \\
    4.2 & Bubble chamber                & Sect.~\ref{subsec:bubble-chamber} \\
    4.3 & $C$, $G$, $P$ at work         & Sect.~\ref{subsec:qnumbers},
                                          Sect.~\ref{subsubsec:bose-symmetry} \\
    4.4 & Ideal mixing                  & Sect.~\ref{subsec:nonets} \\
    \hline
    5.1 & Meson wavefunctions           & Sect.~\ref{subsec:qnumbers} \\
    5.2 & Baryon wavefunctions          & Sect.~\ref{subsec:baryon-antisymmetry},
                                          Sect.~\ref{subsec:su6-56plet} \\
    5.3 & $N$ and $\Delta$ spectrum     & Sect.~\ref{subsec:supermultiplets},
                                          Table~\ref{tab:70-1minus} \\
    5.4 & $J/\psi$ width and $\alpha_s$ & Sect.~\ref{subsec:ozi-width},
                                          Sect.~\ref{subsec:resonances} \\
    \hline
    6.1 & Breit--Wigner, Argand         & Sect.~\ref{subsec:resonances} \\
    6.2 & Quarkonium widths, $\alpha_s$ & Sect.~\ref{subsec:qq-potential},
                                          Sect.~\ref{subsec:ozi-width} \\
    6.3 & The missing $c\bar c$ state   & Sect.~\ref{subsec:charmonium-probes} \\
    6.4 & Properties of the $J/\psi$    & Sect.~\ref{subsec:jpsi-discovery},
                                          Sect.~\ref{subsec:ozi} \\
    \hline
    7.1 & Constituent quark masses      & Sect.~\ref{subsec:magnetic-moments} \\
    7.2 & $\bar p N$ annihilation       & Sect.~\ref{subsec:annihilation},
                                          Sect.~\ref{subsubsec:bose-symmetry} \\
    7.3 & Dalitz plot from data         & Sect.~\ref{subsec:dalitz} \\
    \hline
    8.1 & Dalitz plot                   & Sect.~\ref{subsubsec:dalitz-kinematics} \\
    8.2 & Three-body decays             & Sect.~\ref{subsubsec:dalitz-kinematics} \\
    8.3 & Exotic mesons                 & Sect.~\ref{subsec:gluonic_matter} \\
    \hline
    9.1 & CMS momenta at the boundary   & Sect.~\ref{subsubsec:twobody},
                                          Sect.~\ref{subsubsec:dalitz-kinematics} \\
    9.2 & Baryon naming scheme          & Sect.~\ref{subsec:baryon-naming} \\
    9.3 & Isospin in weak decays        & Sect.~\ref{subsec:deltaI-half} \\
    \hline
    10.1 & Spectroscopic notation       & Table~\ref{tab:l2i2j} \\
    10.2 & Vector meson dominance       & Sect.~\ref{subsec:vmd} \\
    10.3 & Double polarization          & Sect.~\ref{subsec:double-polarization} \\
    10.4 & $\gamma\gamma$ decays        & Sect.~\ref{subsec:qnumbers} \\
    \hline
    11.1 & Dalitz plots, $\sqrt{s}$     & Sect.~\ref{subsubsec:thresholds} \\
    11.2 & Form factors                 & Sect.~\ref{subsec:form-factors} \\
    11.3 & Properties of the nucleon    & Sect.~\ref{subsec:rosenbluth},
                                          Sect.~\ref{subsec:neutron-ff} \\
    \hline
    12.1 & Deep inelastic scattering    & Sect.~\ref{subsec:inclusive-kinematics},
                                          Sect.~\ref{subsec:parton-model} \\
    12.2 & Discovery of $W$ and $Z$     & Sect.~\ref{subsubsec:wz-discovery} \\
    12.3 & Callan--Gross relation       & Eq.~\eqref{eq:callan-gross} \\
    12.4 & Neutrino nucleon DIS         & Eq.~\eqref{eq:nu-beam-production} ff. \\
    \hline
    13.1 & Scaling violation            & Sect.~\ref{subsec:dglap} \\
    13.2 & The $\eta$ meson, thresholds & Sect.~\ref{subsec:eta-gateway},
                                          Sect.~\ref{subsubsec:thresholds} \\
    13.3 & Pentaquarks                  & Sect.~\ref{subsec:insight-outlook} \\
    13.4 & Baryon wavefunction symmetry & Sect.~\ref{subsec:baryon-antisymmetry},
                                          Sect.~\ref{subsec:supermultiplets} \\
    \hline
\end{longtable}
